

% !TeX program = pdflatex
% !BIB program = biber

\documentclass[11pt]{article}

\usepackage[a4paper,margin=1in]{geometry}
\usepackage[T1]{fontenc}
\usepackage[utf8]{inputenc}
\IfFileExists{lmodern.sty}{\usepackage{lmodern}}{}
\usepackage{amsmath,amssymb,amsthm,bm,mathtools}
\usepackage{enumitem}
\usepackage{booktabs}
\usepackage{longtable}
\usepackage{array}
\newcolumntype{P}[1]{>{\raggedright\arraybackslash}p{#1}}
\usepackage[expansion=false]{microtype}
\setlength{\emergencystretch}{3em}
\usepackage{xcolor}
\usepackage{hyperref}
\hypersetup{
  colorlinks=true,
  linkcolor=blue!50!black,
  citecolor=blue!50!black,
  urlcolor=blue!50!black,
  pdftitle={Weak-field gravity in C-Motion: scalar failure and conditional tensor recovery},
  pdfauthor={Toma Kavonius}
}

\usepackage[backend=biber,style=numeric,sorting=none,maxnames=6,uniquename=false,url=false]{biblatex}
\addbibresource{cmotion_foundational_general.bib}
\addbibresource{cmotion_primary_research.bib}
\addbibresource{cmotion_internal_manuscripts.bib}

\AtEveryBibitem{\clearfield{note}\clearfield{annotation}\clearfield{keywords}\clearfield{file}}


\newcommand{\cmotion}{\mbox{C-Motion}}
\newcommand{\ElEmEr}{\mbox{ElEmEr}}
\newcommand{\cmc}{c}


\newcommand{\statusNegative}{\textcolor{red!65!black}{\textbf{Negative}}}
\newcommand{\statusStructural}{\textcolor{blue!55!black}{\textbf{Structural}}}
\newcommand{\statusConditional}{\textcolor{yellow!60!black}{\textbf{Conditional}}}
\newcommand{\statusAssumed}{\textcolor{orange!80!black}{\textbf{Assumed}}}
\newcommand{\statusOpen}{\textcolor{red!75!black}{\textbf{Open}}}
\newcommand{\statusConjectural}{\textcolor{magenta!65!black}{\textbf{Conjectural}}}
\newcommand{\statusDeferred}{\textcolor{gray!70!black}{\textbf{Deferred}}}
\newcommand{\statusModelDerived}{\textcolor{cyan!55!black}{\textbf{Model-derived}}}


\newcommand{\NOne}{\textbf{N1}}
\newcommand{\NTwo}{\textbf{N2}}
\newcommand{\Control}{\textbf{C}}
\newcommand{\OpenBridge}{\textbf{O}}
\newcommand{\RenameRisk}{\textbf{R}}

\theoremstyle{plain}
\newtheorem{proposition}{Proposition}
\theoremstyle{remark}
\newtheorem*{remark}{Remark}

\title{\textbf{Weak-field gravity in \cmotion{}:\\scalar failure and conditional tensor recovery}}
\author{Toma Kavonius}
\date{24 November 2026}

\begin{document}
\maketitle

\begin{center}
\small\itshape
Scope and limitations.  This manuscript tests weak-field gravitational structures under
explicitly stated C-Motion antecedents.  The scalar/rate-only route is excluded within
its stated model, while the successful rank-2/helicity-2 route remains Conditional on
vacuum, source, universality, coupling, and nonlinear-completion assumptions that are
not yet derived from native ElEmEr dynamics.
\end{center}
\vspace{0.6em}

\begin{abstract}
This paper investigates whether weak-field gravity can emerge from an
underlying field/state in physical three-space, with locally $c$-limited
change and recurrent closures. The metric formulation is used as the successful effective
benchmark, not as a prescribed microscopic ontology. We test a minimal
scalar/rate-only response against a conditional massless rank-2 completion.
Under the stated optical assumptions, pure scalar/rate loading yields
$\gamma=0$ and half the observed leading light deflection, despite capturing
aspects of Newtonian response and gravitational clock shifts. This is a
restricted negative result, not a no-go theorem for all scalar extensions.
With independently supplied vacuum, conserved-source, universal-coupling,
common probe-response and linear gauge antecedents, the standard rank-2
completion gives $\gamma=1$, leading light bending, Shapiro delay and two
luminal transverse tensor polarizations. The GR value $\beta=1$ and standard
first post-Newtonian behavior further require the specified nonlinear
completion. Multimessenger GW170817 observations constrain propagation speed
\cite{Abbott2017GW170817Gamma}. These are conditional compatibility results, not a native derivation
of general relativity. We examine a candidate interpretation in which
source-dependent field/stress organization redirects open propagation and
changes recurrent-return conditions, and we calculate bounded closure-projection
diagnostics that test proposed scalar-response alternatives. Native mode
selection, gravitational sources and coupling, finite-closure universality,
field self-stress, the strength $G$, and strong-field dynamics remain open.
Acceleration-induced closure instability is a separate conjecture requiring a
stability calculation.
\end{abstract}


%======================================================================
\section{Introduction: what gravity asks of a field-first theory}
\label{sec:introduction}

General relativity provides the benchmark that any alternative microscopic account of
gravity must meet; weak-field and post-Newtonian tests provide a compact empirical control framework \cite{Will2014Confrontation}.  In its standard formulation, gravitational phenomena are represented
geometrically by a metric and its curvature.  \cmotion{} asks a different microscopic
question: if the underlying physical picture is instead a dynamical field/state in
physical three-space, what must that field be able to do so that the same tested
weak-field phenomena emerge?

The distinction is ontological rather than observational.  A successful \cmotion{}
gravity sector may admit the usual metric description at the effective level.  Nothing
in this paper requires the standard geometric mathematics to be discarded.  The narrower
claim is that such mathematics need not be taken as the primitive material ontology of
the underlying theory.

The shared \cmotion{} framework begins from physical three-space, an underlying dynamical
field/state called \ElEmEr{}, and primitive physical changes that propagate locally at the
universal speed $c$.  Localized recurrent organizations of that field are called
\emph{closures}.  Primitive succession orders physical change, while ordinary time is
operationally read from recurrent physical processes and signalling.  Thus when standard
gravity notation uses a four-coordinate interval or a metric, this paper treats that
structure as an effective/operational comparison language.  Primitive succession is not
being promoted to a fourth spatial dimension.

The gravity problem can then be stated simply:
\begin{quote}
\emph{Can a local-$c$, field-first theory in physical three-space reproduce the observed
weak-field behavior of gravity, and what additional gravitational structure is required
for that to happen?}
\end{quote}

The answer developed below has two parts.  First, the simplest scalar/rate-only route
fails a decisive relativistic test: light bending.  Second, a massless rank-2 completion
with universal stress coupling reproduces the required weak-field tensor structure, but
only under explicit antecedents whose microscopic \ElEmEr{} origin remains open.

Alongside that mathematical argument, the paper keeps one candidate physical reading in
view.  A source-dependent gravitational environment may reorganize the local conditions
under which open propagation is redirected and recurrent closures complete their return
cycles.  The resulting trajectories and clock relations can then acquire an effective
geometry even though primitive propagation remains local at $c$ in physical three-space.
This is an interpretation of the conditional weak-field structure, not an additional
derivation of it.

%======================================================================
\section{What the gravitational sector must reproduce}
\label{sec:targets}

Before choosing a field representation, it is useful to separate the observational target
from the microscopic explanation.  In the weak-field regime relevant here, a viable
gravity sector must reproduce at least the following structures within the scope of this
paper:
\begin{itemize}[leftmargin=*]
  \item Newtonian attraction for slowly moving bodies;
  \item an approximately composition-independent free-fall response, tested
        for specified materials by MICROSCOPE \cite{Touboul2022MICROSCOPE};
  \item the full leading deflection of light and the corresponding Shapiro time
        delay; Cassini radio-link timing constrains their common PPN coefficient
        $1+\gamma$ through the delay \cite{Bertotti2003Cassini};
  \item weak-field clock-rate behavior compatible with gravitational redshift;
  \item the standard parametrized post-Newtonian values $\gamma=1$ and $\beta=1$ where
        the necessary linear and nonlinear antecedents are supplied; lunar laser
        ranging constrains equivalence-principle/PPN combinations, including
        $4\beta-\gamma-3$ in the stated conservative PPN class, rather than
        measuring $\beta$ independently \cite{Williams2004LLR};
  \item gravitational disturbances propagating at $c$ in the admitted massless branch;
  \item two physical transverse tensor polarizations of those disturbances.
\end{itemize}

These are not independent \cmotion{} discoveries.  They are empirical or standard-GR
benchmarks.  None of these measurements identifies a native C-Motion
source or mediator, proves exact equality of a PPN coefficient to unity,
or establishes universal microscopic closure response.  Their value here is discriminatory: they tell us which candidate
realizations of gravity are too simple and which effective field structures remain
compatible with observation.

A few terms will recur.  A \emph{scalar field} assigns one number to each point.  A
\emph{tensor field} carries several linked components; a rank-2 tensor has two indices and
can therefore encode directional and stress-like information unavailable to one scalar.
A \emph{massless mediator} is a propagating field mode without a mass gap in the admitted
effective description, allowing long-range propagation with the relevant relativistic
characteristic structure.  \emph{Universal coupling} means that the gravitational field
couples in the same structural way to the appropriate total source rather than to an
arbitrary composition-specific label.

%======================================================================
\section{C-Motion constraints before choosing a gravity field}
\label{sec:constraints}

\subsection{Local-\texorpdfstring{$c$}{c} causal propagation}

The first inherited constraint is kinematic.  When a physical field/state change is
generated locally and propagates openly, the shared \cmotion{} framework requires that
change to extend along successor-directed local-$c$ characteristics.  Any genuine open
gravity influence realized in the same underlying field must therefore respect the same
local propagation speed.

This condition says nothing by itself about the identity of the gravity mode.  In
particular,
\begin{equation}
\boxed{
\text{primitive local }c
\;\not\Rightarrow\;
\text{massless rank-2 gravity or GR}.}
\label{eq:localc-not-rank2}
\end{equation}
It constrains propagation after a gravitational response exists; it does not derive the
source, rank, helicity, coupling, universality, or field equation.

\subsection{Local source, field reorganization, and compound response}

The second inherited constraint is dynamical in form but not yet sector-specific.  The
shared candidate interaction architecture is
\begin{equation}
\boxed{
\begin{gathered}
\text{gravity-sector source}
\longrightarrow\text{local field/stress reorganization}\\
\longrightarrow\text{steering or compound response}
\end{gathered}}
\label{eq:response-template}
\end{equation}
This is useful because it identifies the physical jobs a gravity theory must perform.  A
source must alter the local complete field state; the altered field/stress structure must
then change the motion or internal organization of an open mode or finite closure; and
any newly generated change must propagate causally at local $c$.

The Main Framework supplies this response template only as a bounded
architectural guide \cite{CmotionMainFramework2026}, not as a gravity law.
It does not identify a gravity source or an overlap of configurations as gravity,
and does not turn a non-null invariant into an activation law.  Overlap is a
geometric/state relation; non-nullness is an invariant property in a specified
field representation.  Cross terms of overlapping individually null components
can be non-null in such a model, but this is not true of every overlap.
Departure from a locally valid free-propagation class provides an interaction
opportunity, while actual dynamical activation still requires its sector's
constitutive response and physical conditions.  These distinctions do not
supply $G$, select a rank-2 mediator, or derive universal coupling.
It describes the shape of the missing bridge, not its solution.

\subsection{The vacuum antecedent}
\label{sec:vacuum}

The conditional construction below uses F to denote the genuinely
emergent/background-independent vacuum branch, with Lorentz-invariant effective
local propagation.  Absence of an observable preferred inertial direction or
velocity is an operational requirement of the tested effective Lorentz sector;
it does not by itself establish F.  A fixed, locally undetectable
Lorentz-compatible background can share that same local propagation law.

The choice between these branches is not derived from local $c$.  F therefore
remains an explicit Open antecedent of the conditional gravity route.
Operational Lorentz compatibility, genuine background independence, native
gauge redundancy, and the successful weak-field calculation are distinct
questions; the latter does not settle the former.
In particular, the standard effective Fierz--Pauli calculation can be
formulated on a fixed Minkowski background.  F specifies the proposed
C-Motion realization selected for investigation, not an additional
background-independence requirement of that mathematical control.
Assuming F also supplies no native proof of M, S, U, P, L or N below.

%======================================================================
\section{The simplest route: scalar/rate loading}
\label{sec:scalar}

The most economical possibility is to let gravity alter only one scalar quantity: the
local propagation or clocking rate.  This is physically attractive because it tries to
connect two familiar weak-field effects with one local environmental factor: slowly
moving matter responds to the gravitational potential, while recurrent physical clocks
change rate relative to clocks elsewhere.  The candidate therefore deserves to be tested
before richer field structure is introduced.  Light provides the sharper test.

\subsection{Exact scope of the tested model}
\label{sec:scalar-scope}

The calculation assumes a static dimensionless scalar load $U(\mathbf x)$,
normalized in the isolated weak-field Newtonian comparison as
\begin{equation}
U(\mathbf x)=-\frac{\Phi_N(\mathbf x)}{\cmc^2},
\qquad \Phi_N=-\frac{GM}{r}
\quad\text{outside a spherical source},
\label{eq:scalar-U-Newton}
\end{equation}
with $U>0$ for the attractive potential in that exterior convention.
Here $\Phi_N$ is the Newtonian potential per unit mass, vanishing at spatial
infinity, and $U$ is dimensionless.  Throughout the observational comparison,
$G>0$ is the externally measured Newton constant.  It normalizes an imported
effective control; neither its attractive sign nor its value is derived from
native ElEmEr dynamics.  Native source, coupling and normalization remain Open.
This is a measured comparison constant under the framework-wide one-scale
policy, not a second independently fitted native sector scale.
No gravity-sector calibration of $\Lambda_C$ is performed here; a future
native completion must account for its coupling within the shared scale policy
rather than relabel an independently fitted parameter as a comparison input.
The
model lets this load act only through one local rate factor
\begin{equation}
N=e^{-U}.
\label{eq:Nscalar}
\end{equation}
The corresponding open-mode propagation speed and optical index are
\begin{equation}
v=\cmc N,
\qquad
n=\frac{\cmc}{v}=e^U.
\label{eq:scalar-index}
\end{equation}
Here $v=\cmc N$ denotes an \emph{effective coordinate/optical speed} in the
chosen static weak-field chart, not the locally measured primitive speed of an
open \ElEmEr{} propagation segment, which remains $\cmc$.  The associated
optical index is likewise a coordinate comparison device, not a primitive
constitutive refractive medium.  No independent spatial or routing tensor
response is added.  The result is read only in the weak-field
geometric-optics/post-Newtonian regime.  It is therefore a test of this
pure scalar/rate-loading class, not a theorem against every scalar--tensor, multichannel,
nonlocal, or strongly nonlinear gravity theory.

\subsection{What the scalar control captures}

Within this restricted weak-field control, the same position-dependent factor can play
two useful roles.  It can encode a Newtonian-like response for slowly moving matter, and
it can change the recurrence rate of a physical clock or closure cycle.  In \cmotion{}
language this makes the scalar picture an attractive first approximation: primitive
segments still propagate locally at $c$, while the organization of a recurrent process
can respond to its environment.

This is not yet a microscopic derivation of either universal free fall or gravitational
redshift.  It says only that a rate-like gravitational sector can account for part of the
observed weak-field behavior without changing the primitive local propagation speed.  The
light-bending calculation tests whether that one scalar is enough.

\subsection{Why light bending exposes the failure}
\label{sec:scalar-failure}

For $U\ll1$,
\begin{equation}
n=e^U=1+U+O(U^2).
\end{equation}
In the standard \emph{isotropic spatial-coordinate} weak-field PPN convention,
we use metric signature $(-+++)$ and $x^0=ct$.  This is a declared convention
change from the Main Framework's $(+---)$ stress convention, not a physical
disagreement.  With $U=-\Phi_N/c^2$ as defined above, the effective metric has $g_{00}=-(1-2U+O(U^2))$ and
$g_{ij}=(1+2\gamma U+O(U^2))\delta_{ij}$, the comparison optical form is
\begin{equation}
n_{\rm PPN}=1+(1+\gamma)U+O(U^2),
\label{eq:ppn-index}
\end{equation}
The index follows from the static null condition as
$n=\sqrt{(1+2\gamma U+\cdots)/(1-2U+\cdots)}$.
Here $\gamma$ is the parametrized post-Newtonian coefficient that measures the spatial
part of the gravitational response relative to the Newtonian potential.  Matching the
pure rate-loaded model to Eq.~\eqref{eq:ppn-index} gives
\begin{equation}
\boxed{\gamma=0.}
\label{eq:gamma0}
\end{equation}
The corresponding leading deflection of a light ray with impact parameter $b$ is
\begin{equation}
\boxed{
\Delta\theta_{\rm scalar}
=\frac{2GM}{b\cmc^2},}
\label{eq:half-bending}
\end{equation}
only half of the $\gamma=1$ value.

\begin{proposition}[Scoped scalar/rate-loading no-go]
Within the model class of Section~\ref{sec:scalar-scope}, pure scalar/rate loading implies
$\gamma=0$ and therefore insufficient leading light bending.
\end{proposition}

The physical lesson is more important than the label.  Correct Newtonian attraction for
slow bodies does not guarantee correct relativistic gravity.  In the standard weak-field
PPN bookkeeping, the rate/temporal contribution supplies only part of the light
deflection; the coefficient $\gamma$ measures the additional spatial/tensor response.
The tested scalar supplies the first part but no independent spatial response, leaving
the familiar ``half-bending'' result.  This ``two halves'' language is only a weak-field
PPN intuition, not a decomposition theorem for arbitrary gravity theories.

The observed bending of light therefore requires additional independent spatial/directional
response information that the tested single scalar does not contain:
\begin{equation}
\boxed{
\text{correct Newtonian attraction}
\;\not\Rightarrow\;
\text{correct relativistic gravity}.}
\end{equation}
This statement deliberately stops at information content.  It does not by itself identify a
shear or trace-free channel, and it does not uniquely select a rank-2 representation.  Those
are possible realizations of the missing information, not consequences of the scalar no-go.

The same scalar control would also give only one third of the standard perihelion
coefficient if one inserts $\gamma=0$ and $\beta=1$ into
\begin{equation}
R_{\rm peri}=\frac{2+2\gamma-\beta}{3}.
\end{equation}
That is a downstream consistency check on the same missing spatial response, not an
independent \cmotion{} perihelion prediction.

%======================================================================
\section{What the scalar failure tells us}
\label{sec:lesson}

The negative result narrows the gravity problem without solving it.  A viable \cmotion{}
gravity sector must carry more spatial or relational information than a single rate
scalar.  That additional structure must affect the propagation of light and must be able
to respond to more than a scalar rest-mass label.

This conclusion does not uniquely prove that gravity must be rank 2.  It tells us why
the \emph{tested amount/rate-only coupling} is insufficient.  A scalar field in a more
general theory may have gradients, derivative couplings, extra source data, or
nontrivial finite-closure projection rules; the negative result does not exclude
those possibilities.  A separate restricted source-capacity test assumes a
rotationally scalar, linear, first-order source depending only on the local
amount $U_0$ (and any already specified flux), without independent angular
stress data, derivative reconstruction or nonlocal information.  It fails to distinguish an equal-budget isotropic radiation bath
from oppositely directed radiation streams: their energy-like scalar totals can
agree, with both fluxes zero, while their second directional moments differ.
Under these assumptions a scalar linear functional of the angular distribution
has rotation-invariant weight and therefore accesses only its total amount.
It cannot distinguish these equal-input configurations.  That is a separate
source-capacity restriction on the \emph{amount-only} model, not a no-go theorem
for all scalar-mediated models.  A symmetric rank-2 field is the standard effective
completion tested here because it can accommodate the additional direction-sensitive
stress/source data while coupling to a relativistic source.

The scalar failure therefore motivates the tensor test without making it a native
uniqueness theorem.

A separate negative result is also important for the broader closure framework.  Gravity
with the observed Newton coupling $G$ is far too weak at elementary-particle scales to be
the universal microscopic binder that maintains an elemental closure.  Thus
\begin{equation}
\boxed{
\text{observed }G\text{-strength gravity}
\neq
\text{universal elementary-closure binder}.}
\label{eq:not-binder}
\end{equation}
This does not prevent gravity from acting on already formed closures.  It only excludes
ordinary observed gravity from doing a different microscopic job for which it is far too
weak.  The actual closure binder remains outside the present Gravity paper.

\subsection{A tested closure-projection alternative and its limits}
\label{sec:closure-projection}
The scalar/rate exclusion does not directly cover a direction-sensitive
scalar steering law followed by a dynamical projection onto a recurrent
closure's collective motion.  The following self-contained projector tests
examine this alternative without establishing its native dynamics.  Assume
unit local directions $\mathbf n$, a positive normalized cycle measure,
and a weak external response approximately uniform across the closure.
Define only a \emph{dimensionless diagnostic response}
\begin{equation}
 \mathbf r(\mathbf n;\hat{\mathbf a})
 =(I-\mathbf n\mathbf n^{\mathsf T})\hat{\mathbf a},
 \qquad
 C_{ij}=\langle n_i n_j\rangle,\qquad
 \overline{\mathbf r}=(I-C)\hat{\mathbf a}.
 \label{eq:projection-diagnostic}
\end{equation}
The unit direction $\hat{\mathbf a}$ and cycle measure are supplied model
inputs.  Equation \eqref{eq:projection-diagnostic} specifies neither a
physical source nor a dimensionful force/response law.  Its identification
with a physical closure-centre response is an additional control assumption.

For a uniformly sampled planar cycle with unit normal $\mathbf N$,
$C=(I-\mathbf N\mathbf N^{\mathsf T})/2$, and the longitudinal diagnostic
coefficient is
\begin{equation}
 \lambda(\mu)
 =\hat{\mathbf a}\cdot\overline{\mathbf r}
 =\frac{1+\mu^2}{2},
 \qquad \mu=\mathbf N\cdot\hat{\mathbf a}.
 \label{eq:projection-planar-coefficient}
\end{equation}
A ray with $\mathbf n_{\rm ray}\cdot\hat{\mathbf a}=0$ has pointwise
longitudinal coefficient
$\hat{\mathbf a}\cdot(I-\mathbf n_{\rm ray}\mathbf n_{\rm ray}^{\mathsf T})
\hat{\mathbf a}=1$ in this projector control, whereas a closure locked
to $\mu=0$ has coefficient $1/2$.
\emph{If} one additionally supplies a common physical coefficient,
compatible operational maps from cycle response to matter acceleration
and ray steering, and a universal half-response lock, calibration on slow
matter leaves a factor-two comparison of these pointwise geometric
coefficients.  It is not a complete integrated bending calculation:
the ray direction and its transverse relation to the field vary along
the trajectory.  A physical deflection angle requires a specified
response law integrated along that trajectory; the matter comparison
also requires a centre-of-energy evolution equation.  This remains a
narrow Conditional diagnostic.
It does not derive the lock, the source, optical-path or clock laws,
$\gamma$ in a complete PPN theory, tensor radiation, or a universal
equivalence principle.  It also does not remedy the independent amount-only
source-capacity restriction: response projection and source law are distinct.

Four narrowing tests follow:
\begin{enumerate}[leftmargin=*]
 \item \textbf{Raw isotropic averaging.}
 Since $\operatorname{tr}C=1$, a universal isotropic raw response
 $I-C=\lambda I$ obeys $3\lambda=2$.  Thus
 \begin{equation}
  \boxed{\lambda_{\rm raw,iso}=\frac23\ne\frac12.}
  \label{eq:projection-two-thirds}
 \end{equation}
 This is exact algebra for the supplied projector and normalized cycle
 measure, not a theorem about every physical closure response.
 \item \textbf{Minimum differential steering strain.}
 In the tested planar diagnostic, the simplest strain criterion prefers
 $|\mu|=1$, the full-response orientation, rather than the half-response
 orientation $\mu=0$.  Appendix~\ref{app:projection-strain} defines
 the variance criterion and evaluates it explicitly.
 \item \textbf{Finite-size tides.}
 In the symmetric circular counterflow diagnostic, first-gradient forcing
 leaves internal tidal strain while its first-order collective average
 cancels.  It does not supply the missing half-response coefficient.
 The spatial expansion and opposite-point symmetry are stated in
 Appendix~\ref{app:projection-tidal}.
 \item \textbf{Retarded first-hit timing.}
 A settled static field has no continuing changing-field first-hit term.
 For changing fields, the tested leading odd timing/memory correction
 produces internal phase strain but no first-order symmetric collective
 coefficient.  Appendix~\ref{app:projection-retarded} gives the timing
 model, expansion and cancellation conditions.
\end{enumerate}

These are scoped Structural/Conditional diagnostics and Negative tests of
particular rescue mechanisms.  A surviving scalar route would require a
native non-raw collective projection, closure-memory law, different local
response, or additional directional sector, together with the other gravity
obligations.  The qualitative cycle-projection interpretation in
Section~\ref{sec:budget} remains Conjectural; recurrence alone does not earn
a universal half-response or universal free fall.

%======================================================================
\section{A conditional massless rank-2 completion}
\label{sec:rank2}

Once the pure scalar route has failed, the constructive question is whether a richer
relativistic field can meet the observational target while remaining compatible with the
shared \cmotion{} constraints.  The completion studied here is a massless symmetric
rank-2 field, using the standard massless kinetic limit of the Fierz--Pauli Lagrangian
(linearized Einstein--Hilbert), as an effective control class
\cite{FierzPauli1939}.

``Rank 2'' means that the effective field carries two indices, $B_{\mu\nu}$.  This gives
it enough component structure to represent a coupled response of energy density,
momentum density, stresses, and directional relations rather than only one scalar
potential.  ``Massless'' means that the admitted weak-field gravitational mode has no
mass gap.  A mass gap would in general modify the long-range behavior and the wave
dispersion, so masslessness is part of the tested long-range completion rather than a
consequence of primitive local $c$.

At the present interpretive level, an unscreened massless gravitational mode fits
naturally on the \emph{open-propagation} side of the \cmotion{} open/closed taxonomy,
whereas massive matter is represented by recurrent closed organization.  This is a
useful ontology picture, not a theorem that every open mode is massless, every massless
mode is open, or gravity is natively selected by that taxonomy.

\subsection{What is assumed before the recovery works}
\label{sec:antecedents}

The conditional chain is easiest to read if its ingredients are separated explicitly.
\begin{longtable}{@{}P{1.2cm}P{4.3cm}P{3.0cm}P{5.0cm}@{}}
\toprule
\textbf{Label} & \textbf{Antecedent} & \textbf{Status} & \textbf{Role} \\
\midrule
\endhead
F & Frameless Lorentz-invariant vacuum / emergent-background branch & \statusOpen{} & Selected C-Motion realization; not an independent requirement for the fixed-Minkowski effective Fierz--Pauli control. \\
M & Massless symmetric rank-2/helicity-2 gravitational mediator & \statusConditional{} & The effective tensor field class whose weak radiative sector is tested; native selection remains Open. \\
S & Suitable symmetric conserved total matter--radiation--interaction source, supplied by an action or equivalent native conservation/source structure & \statusConditional{} & A conserved relativistic source $\mathcal T^{\mu\nu}$ is supplied for the rank-2 control; native source structure remains Open. \\
U & Universal coupling to that total source, with the effective attractive sign and normalization supplied by the external Newtonian comparison using measured $G$ & \statusAssumed{} & Universality is an active antecedent; its native derivation, sign, source identity, and normalization remain Open. \\
P & Common effective metric/probe response for matter, clocks and light & \statusConditional{} & Standard minimal probe-response control in an isotropic weak-field chart; the native response of finite closures remains Open. \\
L & Local two-derivative massless rank-2 gauge/Fierz--Pauli linear class & \statusConditional{} & The standard ghost-free linear tensor structure and its two helicity-2 modes. \\
N & Nonlinear GR-class self-coupling completion, including consistent gauge/constraint structure and field self-source & \statusOpen{} & A generic nonlinear interaction is insufficient for $\beta=1$; the required GR-class completion is not native-derived. \\
\bottomrule
\end{longtable}

The table is the scientific boundary of the positive result.  Local $c$ does not derive M;
the generic interaction architecture does not derive S, U or the probe-response
premise P; and the scalar failure does not prove that L is the unique microscopic
completion.  The result is instead a precise
compatibility statement: if these antecedents are available, the familiar weak-field
tensor phenomenology follows.

\paragraph{Selected realization versus effective control.}
F specifies the proposed C-Motion route with an emergent vacuum and
background independence; that realization remains Open.  The effective comparison instead
supplies M, S, U, P and L on a Lorentz-compatible background, which may
be fixed Minkowski space, with the stated exclusion of extra relevant
scalar modes.  Standard Fierz--Pauli mathematics neither establishes F
nor independently requires proof of microscopic background independence.
The chains below retain F to identify the selected C-Motion route, not
to assert that F is a necessary premise of every effective spin-2 model.

\paragraph{Probe response is not source universality.}
The effective PPN reading also supplies a common operational metric coupling:
with a suitable normalization $g^{\rm eff}_{\mu\nu}=\eta_{\mu\nu}
+\kappa_G B_{\mu\nu}+O(B^2)$, the comparison dynamics of light, clocks,
and slow matter is governed by that same effective metric.  This is P, distinct
from the source-coupling premise U.  The sign and normalization are fixed by
the Newtonian comparison, not by a native derivation.  In particular,
the finite internal return dynamics of a physical closure has not been proved
to follow this metric universally; P is imported solely for the effective
weak-field/PPN control.

\subsection{The linear tensor structure}

With a symmetric weak field $B_{\mu\nu}$, the standard linear gauge redundancy is
\begin{equation}
B_{\mu\nu}\sim B_{\mu\nu}
+\partial_\mu\xi_\nu+\partial_\nu\xi_\mu.
\label{eq:linear-gauge}
\end{equation}
For a local quadratic two-derivative action, write the explicit coefficient
basis, with $B=\eta^{\mu\nu}B_{\mu\nu}$ and all contractions performed
with $\eta_{\mu\nu}=\operatorname{diag}(-1,1,1,1)$.
Thus $\partial^0=-\partial_0$, and raising both indices of $B_{\mu\nu}$
uses two factors of $\eta$, not an independent Euclidean convention:
\begin{align}
\mathcal L_2={}&a\,\partial_\lambda B_{\mu\nu}
                  \partial^\lambda B^{\mu\nu}
+b\,(\partial_\mu B^{\mu\nu})(\partial^\lambda B_{\lambda\nu})
\notag\\
&+c\,(\partial_\mu B^{\mu\nu})\partial_\nu B
+d\,\partial_\lambda B\,\partial^\lambda B.
\label{eq:FP-quadratic-basis}
\end{align}
Requiring gauge variations to leave the action invariant, up to boundary terms,
fixes the coefficient ratios in this basis:
\begin{equation}
2a+b=0,\qquad b+c=0,\qquad c+2d=0,
\label{eq:FP-kinetic-relations}
\end{equation}
so $(a,b,c,d)\propto(1,-2,2,-1)$, up to an overall normalization and
equivalent integrations by parts.  Reversing the metric signature changes
the overall sign of this kinetic basis, not these coefficient ratios.
For the present signature the healthy TT kinetic term requires $a<0$,
so the coefficient of $(\partial_0 B^{\rm TT}_{ij})^2$ is positive.
The absolute normalization and source-coupling sign are additional inputs,
not predictions of gauge invariance.

For curvature comparisons we take
$R^\rho{}_{\sigma\mu\nu}=\partial_\mu\Gamma^\rho_{\nu\sigma}
-\partial_\nu\Gamma^\rho_{\mu\sigma}
+\Gamma^\rho_{\mu\lambda}\Gamma^\lambda_{\nu\sigma}
-\Gamma^\rho_{\nu\lambda}\Gamma^\lambda_{\mu\sigma}$,
$R_{\sigma\nu}=R^\rho{}_{\sigma\rho\nu}$, and the corresponding Einstein
equation $G_{\mu\nu}=8\pi G\mathcal T_{\mu\nu}/c^4$ in the supplied GR
control.  The PPN and optical formulas above require no imported curvature
sign; the tidal discussion below is qualitative and does not import an
oppositely signed component formula.
The radiative sector then contains two physical helicity-2 modes.  This mathematics is a
standard consequence of the assumed effective field class.  Its role here is to show
that such a class is compatible with the \cmotion{} constraints, not to claim that the
microscopic \ElEmEr{} dynamics has already selected it.

\subsection{Why the source must be more than rest mass}
\label{sec:source}

Relativistic gravity must respond not only to rest-mass density but to the appropriate
total energy--momentum/stress structure.  The effective source is therefore a symmetric
tensor $\mathcal T^{\mu\nu}$ containing energy density, momentum density, momentum flux,
and stresses.  In the $x^0=ct$, $(-+++)$ comparison,
$\mathcal T^{00}$ is energy density,
$\mathcal T^{0i}$ is energy flux divided by $c$, and
$\mathcal T^{ij}$ is spatial momentum flux/stress.  In particular
$\mathcal T=\eta_{\mu\nu}\mathcal T^{\mu\nu}
=-\mathcal T^{00}+\sum_i\mathcal T^{ii}$.
No $(+---)$ trace identity is inserted without this conversion.

If the complete effective system admits one suitable autonomous local Poincar\'e-invariant
action, translation invariance supplies a conserved energy--momentum current and Lorentz
invariance permits its symmetric Belinfante--Rosenfeld form.  Under that standard control
premise,
\begin{equation}
\partial_\mu \mathcal T^{\mu\nu}=0.
\label{eq:source-conservation}
\end{equation}
An action is a strong route to this structure, but the public claim does not require the
stronger assertion that only a variational formulation could ever provide it.  An equivalent
native conservation/source construction would serve the same linear role if it supplied the
required local object consistently.

The conservation law is important both for consistent coupling to the gauge-redundant
rank-2 field and for the radiation multipole structure discussed later.  Physically, this
also explains why ``mass causes gravity'' is too narrow a slogan: radiation, momentum
flow, pressure, and stress belong to the relativistic source accounting as well as rest
energy.  One \cmotion{} reading is to call this the complete budget/stress accounting of
the organized field.  That wording is mnemonic only; it does not derive
$\mathcal T^{\mu\nu}$ from a microscopic propagation budget.

The remaining native problem is sharper than the phrase ``gravity couples to stress-energy.''
\cmotion{} still needs a local source functional: a rule that takes a generic localized
\ElEmEr{} organization and tells the gravitational response which source information it
presents.  Conserved total energy and momentum may be usable as coarse black-box exterior data
when they are independently available, but those totals do not determine the complete local
stress and flux structure.  The native source functional, its conservation law, and its map
into the long-range gravitational response remain open.

\subsection{Universal coupling}
\label{sec:universal}

Universal coupling means that the same gravitational field responds to the appropriate
total source in the same leading way for different forms of matter.  That is what allows
a gravitational response to be approximately composition independent rather than tied
to an arbitrary particle label.

A candidate physical reading is that if gravity responds to the complete organization of
energy, momentum, and stress rather than to a material-specific label, composition
independence becomes understandable.  That intuition is useful, but it does not prove the
weak equivalence principle.

The Main Framework supplies a useful prior distinction
\cite{CmotionMainFramework2026}.  Given a complete isolated recurrent
organization with conserved total energy and momentum and a physical rest
frame, the further effective Lorentz relation
$E^2-c^2|\mathbf P|^2=M^2c^4$ identifies $Mc^2=E_{\rm rest}$.
Its small-momentum expansion then supplies the usual inertial coefficient
within that Conditional effective description.  Neither step derives the
recurrent solution, establishes its gravitational source identity, or proves
composition-independent free fall.  Rest energy, inertial response, and
gravitational source/response universality are separate obligations.

The present paper requires such coupling in the successful weak-field completion, but it
does not derive why \ElEmEr{} realizes it.  The source identity, sign, universality, and
the normalization measured as $G$ remain open microscopic questions.

A second obligation is logically separate:
\begin{equation}
\boxed{
\text{universal gravitational sourcing}
\;\neq\;
\text{universal response of finite closures}.}
\label{eq:sourcing-vs-response}
\end{equation}
Even if every localized organization excites the same gravitational sector with one coupling,
it does not yet follow that differently structured recurrent closures acquire exactly the same
center-of-energy acceleration in a slowly varying field.  A native equivalence-principle
result therefore needs both a universal source/coupling law and a finite-closure response
statement controlling internal-structure corrections.  Assuming the first and recovering
universal free fall in the effective theory is not a derivation of either native obligation.

%======================================================================
\section{Weak-field recovery}
\label{sec:weakfield}

The positive result of the conditional tensor route is that it reproduces the same
weak-field structure that distinguishes general relativity from the failed pure-scalar
candidate.

\subsection{PPN in plain language: \texorpdfstring{$\gamma$ and $\beta$}{gamma and beta}}

The parametrized post-Newtonian (PPN) framework is a standard way to compare weak-field
gravity theories without committing to their microscopic interpretation.  Two parameters
are especially relevant here.  The parameter $\gamma$ measures the spatial/tensor
response associated with a Newtonian potential.  The parameter $\beta$ measures the
leading nonlinear gravitational self-interaction contribution.

For the admitted linear tensor sector within the selected C-Motion route,
\begin{equation}
\boxed{
\mathrm F+\mathrm M+\mathrm S+\mathrm U+\mathrm P+\mathrm L
+\text{no relevant extra scalar admixture}
\Longrightarrow \gamma=1.}
\label{eq:gamma1}
\end{equation}
F labels the route selection in this sufficient chain; the standard
effective calculation does not establish its microscopic vacuum premise.
This is the universally coupled, common-metric massless spin-2
weak-field control; it is not a claim that sourcing alone determines how
\ElEmEr{} closures or light probes respond.  The leading light deflection is then
\begin{equation}
\boxed{
\Delta\theta=\frac{4GM}{b\cmc^2},}
\label{eq:full-bending}
\end{equation}
and the standard leading Shapiro delay coefficient follows from the same weak-field
structure.  The contrast with Section~\ref{sec:scalar-failure} is the central comparison
of the paper: the tested scalar/rate candidate gives $\gamma=0$, whereas the admitted
rank-2 class gives $\gamma=1$.

The linear theory does not by itself fix $\beta$.  Here N is the
\emph{specific} nonlinear GR-class completion of the universally coupled
massless spin-2 theory, with appropriate constraint/gauge consistency, the
same physical coupling normalization and no extra relevant gravitational
sector.  A generic nonlinear self-interaction is not enough.  The standard
self-coupling consistency construction \cite{Deser1970SelfInteraction}
is used only within its specified local, gauge-consistent spin-2 and
conserved-source class; it does not establish that native \ElEmEr{}
realizes those premises.  If N holds,
\begin{equation}
\boxed{
\mathrm F+\mathrm M+\mathrm S+\mathrm U+\mathrm P+\mathrm L+\mathrm N_{\rm GR}
\Longrightarrow \beta=1.}
\label{eq:beta1}
\end{equation}
With $\gamma=\beta=1$, the standard first post-Newtonian perihelion coefficient follows.
A future quadratic action, Hamiltonian, or equivalent conserved-energy seed would still not
be enough to establish this nonlinear step.  Such a structure could support linear mode
dynamics, linear source coupling, and a quadratic field-energy/stress definition, whereas
$\beta$ probes how the gravitational field's own source modifies the static field.  The
nonlinear self-source---equivalently cubic or higher-order data in an action formulation---is
therefore a distinct open gate.

Numerical integrations that are seeded with the standard GR/PPN equations are therefore
implementation controls, not independent evidence for a microscopic \cmotion{} gravity
law.

PPN agreement has a precise but limited meaning.  It says that the admitted effective
weak-field sector reproduces the corresponding tested coefficients.  It does not show
that the full nonlinear Einstein equations have been derived from \ElEmEr{}, nor does it
settle strong-field interiors, mergers, cosmology, or the microscopic source law.

\subsection{Clock behavior without treating time as a substance}
\label{sec:clocks}

Weak gravitational fields also change the relative rates of physical clocks.  In the
standard metric description this effect is encoded in the time-time part of the weak
metric and is described as gravitational redshift or time dilation.  In \cmotion{}, the
operational reading is different: clocks are recurrent physical processes, so a
gravity-dependent change in clock rate means that the gravitational environment changes
the recurrence/return rate of one physical clock relative to another.

This does not mean that a fundamental time substance literally flows more slowly in one
place.  A candidate \cmotion{} interpretation is instead that the gravitational
environment changes the conditions under which a recurrent organization completes one
cycle.  The scalar/rate control of Section~\ref{sec:scalar} shows why such a reading is
natural: a local rate factor can reproduce the leading clock-rate tendency while
primitive propagation remains local at $c$.  The same restricted picture is not enough
for light bending, which is exactly where the missing spatial/tensor response appears.

The effective metric description remains a valid comparison language for the observed
clock relation, while the microscopic \ElEmEr{} mechanism that produces the universal
gravitational clock response remains part of the open source-to-response bridge.

%======================================================================
\section{Gravitational waves}
\label{sec:gw}

A successful gravity theory must also support propagating disturbances, not only a static
field around matter.  In the admitted massless rank-2 sector, small disturbances of the
effective gravitational field propagate as tensor waves.

The field-first picture makes the causal story concrete without supplying a new gravity
law.  A changing source first generates a local gravitational response; the generated
field/state disturbance then propagates outward through physical three-space along the
same successor-directed local-$c$ characteristics used for other open changes.  This
generation-versus-propagation distinction is inherited from the shared framework and does
not determine the source equation.

The \cmotion{} framework therefore supplies one independent kinematic restriction: a
genuine open field/state change must propagate locally at $c$.  The identification of the
physical gravity wave as a massless rank-2 mode, however, additionally requires the M/L
antecedents.  Gravitational-wave speed alone does not select rank 2.
The GW170817/GRB 170817A timing analysis in ApJL \textbf{848}, L13
\cite{Abbott2017GW170817Gamma} constrains propagation-speed differences
under its emission-delay and propagation assumptions.  It does not establish
an exact zero difference, a native gravity source, or the polarization count.

\subsection{Transverse-traceless structure and polarization}

A wave is \emph{transverse} when its physical deformation acts in directions perpendicular
to its propagation direction.  In the rank-2 gauge class, residual gauge freedom removes
nonphysical components and leaves a \emph{transverse-traceless} (TT) radiative field.  The
traceless condition removes the scalar-like trace component.  Two independent physical
tensor polarization patterns remain, corresponding to the two helicity-2 states.

Thus, under M+L,
\begin{equation}
\boxed{
\text{massless rank-2 linear sector}
\Longrightarrow
\text{two TT tensor polarizations}.}
\end{equation}
The physical picture in \cmotion{} is a propagating gravitational field/state change in
physical three-space whose weak-field mathematical description is equivalent to the usual
tensor metric perturbation.  The paper does not treat ``ripples in a material spacetime
fabric'' as the primitive \cmotion{} ontology.

\subsection{Why quadrupole radiation is leading}
\label{sec:quadrupole}

Radiation also tests the source structure.  In the admitted massless helicity-2
control, take an isolated, spatially localized source in the weak-field,
slow-motion regime, use appropriate outgoing far-zone boundary conditions,
and assume no additional light radiative scalar/vector charges.  Then source
conservation implies a constant monopole and nonradiating mass dipole at
leading order, so the first ordinary tensor-radiative multipole is quadrupolar:
\begin{equation}
\boxed{\begin{gathered}
\text{conserved, localized slow source}\\
+\ \text{far-zone pure helicity-2 sector}\\
\Longrightarrow\ \text{quadrupole-leading radiation.}
\end{gathered}}
\label{eq:quadrupole-conditional}
\end{equation}
In standard leading-order notation, the trace-free source quadrupole is
$I_{ij}^{\rm TF}=\int\rho\,(x_i x_j-\delta_{ij}r^2/3)\,\mathrm d^3x$,
with radiative amplitude controlled by its second time derivative.
Here $\rho$ is the effective slow-source mass density in a spatial volume
integral, not the direction-resolved energy/budget density used in
Section~\ref{sec:micro-effective-energy}.
These conditions belong to the effective radiation control and are not
consequences of primitive local $c$ or a native closure-source calculation.

%======================================================================
\section{Free fall, support, tides, and the equivalence-principle boundary}
\label{sec:freefall}

In an effective GR description, freely falling bodies follow geodesics---paths that are
locally straight with respect to the effective gravitational geometry.  That is the
successful macroscopic representation.  The microscopic \cmotion{} question is what
field response of an extended recurrent closure could have that geodesic limit.

One candidate physical reading separates three situations.  A freely falling closure
follows the local gravitational response without a supporting force continually opposing
that motion.  A supported or stationary closure does the opposite: some non-gravitational
interaction must continually resist the local free-fall response, and the body therefore
has nonzero proper acceleration.  Neither statement derives the microscopic gravitational
coupling; it only translates the effective distinction into the field-first ontology.

A second useful distinction is between an approximately uniform response across a small
closure and spatial variation of that response across its finite extent.  In a weak-field
potential notation the first is associated schematically with a local gradient,
$\nabla U$, while the leading differential effect involves variation of that gradient,
schematically $\nabla\nabla U$.  Following the local free-fall response can remove the
uniform part from the immediate experience of a small body; it cannot remove a
difference between the near and far sides of an extended body.  That residual is the
weak-field intuition behind tidal response.

This is not a theorem that closure averaging derives universal free fall, and the Hessian
of one scalar potential is not being identified with the full relativistic Riemann
tensor.  Standard GR packages the general tidal residue in curvature/geodesic-deviation
language.  The \cmotion{} microscopic source-to-field map and the finite-closure response
law remain open, so a native equivalence-principle derivation remains open as well.

The distinction nevertheless matters later: \emph{proper acceleration} of a supported
closure and \emph{tidal variation} across a freely falling closure are different physical
stress questions and must not be conflated.

%======================================================================
\section{A candidate physical reading: budget geometry}
\label{sec:budget}

The weak-field mathematics above can be given a specifically \cmotion{} interpretation
without promoting that interpretation into evidence.  We use ``budget geometry''
here as an interpretive term for the proposed organization of propagation and
return conditions.  Primitive propagation remains locally at $c$, while
recurrence, translation and collective response describe the organization of
the field.  This wording is not attributed to the current Main Framework and
does not supply a conserved scalar substance or a native metric law.

At the present interpretive level, a source-dependent gravitational environment changes
the local conditions under which propagation is redirected and recurrent closures
complete their return.  Schematically,
\begin{equation}
\boxed{
\begin{gathered}
\text{source}
\longrightarrow
\text{local field/stress organization}\\
\longrightarrow
\text{changed propagation/return conditions}
\longrightarrow
\text{compound response}
\end{gathered}}
\end{equation}

A particularly economical microscopic realization would keep the local gravitational
response independent of whether the affected ElEmEr propagation belongs to an open mode
or to a recurrent closure.  On this reading, each locally propagating part of the field
responds to the same gravitational organization according to the same local law.  The
difference appears only at the level of the organized response.  An open mode carries
that local response directly into its continuing propagation, whereas a recurrent
closure samples the response through many internally directed segments of its return
cycle.  Its observable centre-of-organization response could therefore contain only the
part that survives projection or cancellation over the complete recurrence:

\[
\boxed{
  \text{same local gravitational response}
  \longrightarrow
  \begin{cases}
    \text{open propagation: direct propagated response},\\[2pt]
    \text{recurrent closure: cycle-projected net response}.
\end{cases}}
\]

This possibility would not require gravity to distinguish microscopically between
\emph{open} and \emph{closed} ElEmEr.  Openness or recurrence would instead determine how one
common local response appears after the field organization is coarse-grained to an
observable trajectory or closure motion.  It could also provide a route by which a
deformation of the recurrent return pattern under proper acceleration changes the
cycle projection before any eventual loss of closure stability.  At present this is
only a candidate realization: no universal physical closure-cycle averaging law,
physical cancellation factor, acceleration dependence, or equivalence-principle
theorem has been derived.  The diagnostic coefficients and failed simple
mechanisms in Section~\ref{sec:closure-projection} constrain this proposal;
they do not supply its native projection law.


In this picture gravity is not a Newtonian force acting on primitive point-particles.
Free fall is the natural collective response to the local gravitational organization;
weight corresponds to support resisting that response; clock shifts reflect altered
recurrence conditions; and tides reflect spatial variation of the response across an
extended body.

This also gives a careful sense in which gravity can be ``geometric'' without requiring
four-dimensional metric geometry to be the primitive material ontology.  If the local
rules governing propagation and recurrent organization vary with source and position,
the resulting trajectories and clock comparisons acquire an effective geometry.  The
standard metric then remains the precise weak-field mathematical description of that
sector.

No independent derivation follows from this language.  A single universal scalar
``budget substance'' has not been established, microscopic \cmotion{} variables have not
been mapped uniquely into metric variables, and the Candidate Interaction--Response
architecture does not supply the gravitational source or field equation.  The paper's
evidential weight remains the scoped scalar exclusion and the explicit conditional
rank-2 completion.  Budget geometry is a candidate physical reading of those results.

%======================================================================
\section{Acceleration-induced closure instability}
\label{sec:acceleration}

The current Gravity scope also contains a conjectural finite-closure stability lead.  It
is logically separate from the weak-field tensor recovery.  The physical idea is simple:
a recurrent closure must continually maintain a self-consistent return organization, so
sufficiently strong \emph{proper acceleration} might stress, de-lock, reorganize, or
destabilize that recurrence.  This is a candidate interpretation, not an established
effect.

Dimensional reasoning identifies a possible scale, not a proof that a
physical instability occurs at any finite acceleration.  Only a native
accelerated-closure stability calculation could establish a threshold.
The candidate scale is
\begin{equation}
\boxed{
a_{\rm break}\sim C_{\rm dyn}\frac{\cmc^2}{R},}
\label{eq:abreak-R}
\end{equation}
where $R$ is the closure size and $C_{\rm dyn}$ is an unknown dimensionless coefficient.
The physical breakdown hypothesis is Conjectural; the encompassing
native instability problem remains Open (sharp).  Promoting this beyond a
conjectural scale requires an accelerated-closure return,
monodromy/Floquet, or leakage-mode stability calculation.

Only if an independent mass--radius relation
\begin{equation}
R=C_R\frac{\hbar}{m\cmc}
\label{eq:mass-radius}
\end{equation}
is independently supplied, with positive $m$ and $C_R$, does
Eq.~\eqref{eq:abreak-R} become the following Conditional algebraic rewriting
\begin{equation}
\boxed{
a_{\rm break}
\sim
\frac{C_{\rm dyn}}{C_R}\frac{m\cmc^3}{\hbar}
\equiv C_a\frac{m\cmc^3}{\hbar}.}
\label{eq:abreak-m}
\end{equation}
Neither $C_R$, $C_{\rm dyn}$, nor $C_a$ is presently derived.  No value $C_a=1$, no bound
$0<C_a<1$, and no Unruh $2\pi$ coefficient is claimed here.
The use of $\hbar$ in the possible mass--radius bridge adopts the empirical
action-scale identification discussed in Quantum Foundations; it neither
derives the action quantum nor supplies a universal nonlinear
energy--frequency product law.

The Conditional label applies to the substitution of the stated premises,
not to an independently established physical acceleration threshold.  It does
not strengthen the Conjectural physical antecedent.

\paragraph{An idealized nonadiabatic mixing control.}
A distinct effective two-state calculation assumes the Hamiltonian
$H(t)=cp(t)\sigma_3+\Delta_{\rm cl}\sigma_1$, a constant positive
energy gap parameter $\Delta_{\rm cl}$ (level separation $2\Delta_{\rm cl}$
at $p=0$), the supplied rule $i\hbar\partial_t\psi=H\psi$, and prescribed
one-dimensional constant-proper-acceleration kinematics
$dp/dt=Ma_{\rm prop}$ with $M=\Delta_{\rm cl}/c^2$.
The parameter $\Delta_{\rm cl}$ has units of energy.  By contrast,
the gap parameter $\Delta$ in the Quantum Foundations effective
operator has units of angular frequency.  If the two descriptions
refer to the same normalized effective gap, their parameters obey
$\Delta_{\rm cl}=J_0\Delta$; this correspondence is not a native
identification of the two models.
Here $t$ is the effective coordinate time of the control, not primitive
succession promoted to a universal clock.  Setting $p(0)=0$ gives the
linear avoided-crossing sweep
$H(t)=(\Delta_{\rm cl}a_{\rm prop}/c)t\sigma_3+\Delta_{\rm cl}\sigma_1$.
For the ideal full sweep $t:-\infty\to+\infty$, the Model-derived
probability of not remaining on the same instantaneous energy branch is
\[
 P_{\rm LZ}=\exp\!\left[-\pi\frac{a_*}{a_{\rm prop}}\right],
 \qquad a_*:=\frac{\Delta_{\rm cl}c}{\hbar}=\frac{Mc^3}{\hbar}.
\]

This is the standard Landau--Zener transition formula
\cite{Landau1932EnergyTransferII,Zener1932NonAdiabatic},
evaluated for the specified two-state Hamiltonian
and idealized full-sweep protocol.

This analytic control needs no mass--radius premise.  It assumes the
Hamiltonian, gap, phase normalization and entire acceleration protocol;
it does not derive the force coupling or an accelerated self-bound closure.
Branch mixing is not physical closure destruction, stability loss, or
a universal finite breakdown threshold.  Finite pulses, half-sweeps and
different response laws can change the transition history.  This result
fixes neither $C_{\rm dyn}$, $C_R$ nor a universal breakdown coefficient.

Proper acceleration and free-fall tides are different physical questions.  Ideal free
fall does not carry the same proper acceleration as a supported body.  Any disruption of
a freely falling finite closure would instead have to arise from spatial variation of
the gravitational response across the closure---a tidal/gradient effect requiring its
own calculation:
\begin{equation}
\boxed{
\text{proper-acceleration breakdown}
\neq
\text{tidal breakdown}.}
\end{equation}

A genuine accelerated-closure instability would also need to be distinguished
from the standard accelerated-detector Unruh effect \cite{Unruh1976}
and Hawking radiation \cite{Hawking1975}.
The present paper neither answers that comparison nor reinterprets these
standard quantum-field results.  A physical comparison requires the missing
accelerated-closure spectrum first.

%======================================================================
\section{What remains to be derived}
\label{sec:open}

The weak-field result is coherent at its stated conditional level.  Its premises
make the remaining native problem specific.  The missing physics is no longer usefully summarized as ``derive the full
\ElEmEr{} action.''  The open bridge can be organized as
\begin{equation}
\boxed{
\begin{array}{c}
\text{physical vacuum perturbation variables}\\[2pt]
\downarrow\\[-1pt]
\text{local response law and supported mode content}\\[2pt]
\downarrow\\[-1pt]
\text{generic localized source functional}\\[2pt]
\downarrow\\[-1pt]
\text{universal source coupling}\\[2pt]
\downarrow\\[-1pt]
\text{universal finite-closure response}\\[2pt]
\downarrow\\[-1pt]
\text{field energy/stress and self-source}\\[2pt]
\downarrow\\[-1pt]
\text{nonlinear completion}
\end{array}}
\label{eq:missing-bridge-hierarchy}
\end{equation}

\subsection{What is actually propagating?}

Primitive local $c$ tells us how an already identified generated disturbance propagates.  It
does not tell us what the disturbance is.  Before a native gravity mode can be classified,
\cmotion{} must identify a physical perturbation of the unexcited \ElEmEr{} state, schematically
\begin{equation}
X=X_0+\delta X,
\end{equation}
and a local response law for that perturbation.  A compact diagnostic notation is
\begin{equation}
\mathcal D_{AB}\,\xi^B=\mathcal J_A,
\label{eq:diagnostic-native-response}
\end{equation}
where $\xi^A$ stands only for as-yet-unidentified physical perturbation variables.  Neither
$\mathcal D$ nor $\mathcal J$ is a proposed final field equation.  The notation exposes what
is missing: without a native local response operator the framework cannot determine the
vacuum spectrum, mass gap, component count, constraints or redundancies, tensor rank,
helicity, or number of physical polarizations.  Local-$c$ propagation constrains the
characteristics only after such a mode has been physically defined.

This also sharpens the scalar result.  The tested scalar/rate variable is insufficient, so an
admissible theory needs additional independent spatial/directional response information.
Shear, trace-free tensor response, or a rank-2 variable are possible mathematical
realizations, but the negative result does not choose among them.

\paragraph{A physical-characteristic test.}
A future loaded response operator should be checked on its physical
polarizations after the relevant constraints and gauge redundancies have
been handled.  Does its characteristic polynomial retain one common
Lorentzian quadratic cone, schematically
\begin{equation}
 P_{\rm phys}(x,k;U)
 \propto \bigl[\mathcal{G}^{\mu\nu}(x;U)k_\mu k_\nu\bigr]^m,
 \label{eq:physical-cone-test}
\end{equation}

Here $\mathcal{G}^{\mu\nu}$ denotes a candidate effective inverse
cone metric, not an independently derived native gravitational metric.

with $m$ the physical multiplicity, or does loading split the mode cones?
An unloaded Maxwell-like control illustrates a common-cone
case, not a native gravity derivation.  The loaded constitutive law and its
physical characteristic structure remain Open.  A single cone would
establish neither a spin-2 degree-of-freedom count nor common matter/clock
coupling P, and would not fix a conformal factor or nonlinear dynamics.

A separate limitation concerns circular matching.  Prescribing the
optical index that yields full bending, or choosing equivalent reciprocal
routing and constitutive coefficients, restates the missing spatial
response unless those coefficients follow from independently specified
dynamics.  Agreement after that choice is a control, not a new derivation.

\subsection{Source, universality, and finite closures}

Once the response variables are known, a second problem is to say what a generic localized
\ElEmEr{} organization presents to them.  Writing a schematic source functional
\begin{equation}
\mathcal J_A[\mathcal O]
\end{equation}
is useful only as a placeholder for that missing constitutive law.  Rest mass is too little
information for relativistic gravity, while coarse total energy and momentum alone do not
supply the full local stress and flux structure.  The native source object, its conservation
law, coupling sign, and normalization observed as $G$ remain open.

The equivalence-principle residue then divides into two independent questions.  First, does
one source law couple all localized organizations universally to the gravitational response?
Second, do different finite recurrent closures respond universally to a slowly varying field?
The second question can require genuine closure dynamics even if the first has been solved.
This is why Eq.~\eqref{eq:sourcing-vs-response} is more informative than simply saying that
``WEP remains open.''

\subsection{Microscopic budget accounting versus effective gravitational energy}
\label{sec:micro-effective-energy}

An underlying local field might admit genuine local energy, momentum, stress
and flux bookkeeping without its effective massless spin-2 description
possessing a local gauge-invariant gravitational stress tensor.  These are
different propositions.  Angular integration gives the following identities,
\emph{conditional on an assigned conserved exterior budget density} $\rho$
and source-free direction-resolved propagation in physical three-space,
\begin{equation}
\partial_s\rho+\cmc N_i\partial_i\rho=0,\qquad|\mathbf N|=1.
\label{eq:conditional-free-transport}
\end{equation}
For $U_0=\int\rho\,\mathrm d\Omega$,
$F_i=\cmc\int\rho N_i\,\mathrm d\Omega$, and
$M_{ij}=\int\rho N_iN_j\,\mathrm d\Omega$, angular integration yields
\begin{equation}
\boxed{\partial_s U_0+\partial_iF_i=0,\qquad
\partial_s F_i+\cmc^2\partial_jM_{ij}=0.}
\label{eq:free-transport-moments}
\end{equation}
Here $\rho=\rho(\mathbf x,\mathbf N,s)$ is a direction-resolved budget/energy
density per solid angle, $U_0$ its scalar amount per volume, $F_i$ its flux,
and $M_{ij}$ its angular stress moment.  They are not four names for the
same scalar information.  Spatial components and angular moments use the
Euclidean metric of physical three-space, not Lorentzian index raising.
Here $s$ is the compatible succession/evolution label, not a universal
observer clock or a fourth spatial-like dimension.  The moment identities
are exact within this \emph{free-transport model}; they do not identify
$\rho$ as the complete interacting \ElEmEr{} energy density, nor do they
derive closure conversion, a gravitational source operator, or conserved
interacting self-stress.

The dynamical role of directional information is explicit on writing
\begin{equation}
 M_{ij}=\frac{U_0}{3}\delta_{ij}+Q_{ij},\qquad
 \partial_sF_i
 =-\frac{\cmc^2}{3}\partial_iU_0-\cmc^2\partial_jQ_{ij}.
 \label{eq:source-anisotropy-moment}
\end{equation}
Here $Q_{ij}$ is the trace-free \emph{local angular second moment},
not the spatial radiating quadrupole $I_{ij}^{\rm TF}$ of
Section~\ref{sec:quadrupole}.  An isotropic radiation bath and balanced
counterstreams can share $U_0$ and $F_i=0$ while differing in $Q_{ij}$.
When its divergence or boundary flux matters, an amount-only description
therefore loses information needed even for this supplied transport model.
The Negative source-capacity conclusion is restricted to the rotationally
scalar, linear, first-order amount-only class stated in Section~\ref{sec:lesson}.
It does not forbid nonlinear, derivative, nonlocal or additional-state scalar
models, and does not select a native rank-2 source.  Identification of these
free-transport moments with a complete interacting gravitational source
remains Open.

At the \emph{effective rank-2} level, ordinary GR-class gravitational
self-energy is not generally an independent local gauge-invariant,
Lorentz-covariant tensor of the same kind as the matter source.  Gauge-
dependent stress descriptions, fluxes and boundary charges can still carry
physical energy information.  This effective limitation does \emph{not}
prohibit microscopic ElEmEr local energy/stress, but a projection from one
level to the other has not been constructed.

There is a further precise compatibility test: the Weinberg--Witten
energy--momentum obstruction \cite{WeinbergWitten1980} concerns massless
helicity-2 states and an
appropriate conserved Lorentz-covariant stress tensor under its quantum
field-theoretic hypotheses.  The standard effective gauge/GR description
avoids this obstruction through its gravitational energy representation;
that does not establish the same escape for a prospective microscopic
ElEmEr theory.  The native Lorentz symmetry, physical state space,
energy--momentum operator and relevant matrix elements are not derived.
Consequently neither an unavoidable contradiction nor a native evasion of
the theorem has been demonstrated.  In particular, a posited frameless
vacuum or the word ``emergent'' is not on its own an evasion proof.

\subsection{The nonlinear stress/self-coupling gate}

A linear response law does not by itself say how much energy, momentum, or stress belongs to
the gravitational disturbance.  Nonlinear self-coupling therefore needs a consistent native energy,
source and self-response structure.  It need not make the \emph{effective}
gravity self-energy an ordinary local gauge-invariant tensor; see
Section~\ref{sec:micro-effective-energy}.  Schematically, the target would have the form
\begin{equation}
\mathcal D\xi
=\kappa_G\Big(\mathcal J_{\rm matter}+\tau[\xi]\Big),
\label{eq:diagnostic-self-source}
\end{equation}
where $\tau[\xi]$ denotes the gravitational field's own source contribution.  Equation
\eqref{eq:diagnostic-self-source} is a completion target, not a derived \cmotion{} equation:
$\kappa_G$, $\mathcal J_{\rm matter}$, and $\tau$ are all unresolved natively.

An arbitrary operator $\mathcal D$ does not uniquely define $\tau[\xi]$.  One possible route
is an action-integrable response with a native pairing; another is a Hamiltonian/symplectic
construction; an equivalent conserved stress-functional structure could in principle serve
the same physical role.  The important requirement is not the word ``action'' but a
consistent definition of field energy/stress from the same native dynamics that defines the
response and source, together with its compatibility with the effective
gauge/constraint representation.  If a variational or Hamiltonian route is pursued, the next targets are a
physical perturbation variable, a native pairing or symplectic structure, and a compatible
response operator.  None has yet been obtained.

Three consistency tests make this target more concrete: the matter source
and field self-source must belong to a compatible source/conservation and
transformation structure; their coupling must use the same physical
normalization; and nonlinear evolution must preserve the constraints that
removed unwanted propagating modes in the admitted linear sector.
These are obligations, not tests already passed by native gravity.

A separate selection issue remains.  Within reduced constitutive controls,
the stated kinematic, conservation, symmetry and consistency constraints
admit inequivalent laws; admitting a class of laws is not yet physical selection.
Merely replacing an arbitrary response operator by an arbitrary action,
Hamiltonian, stress functional, or memory kernel can relocate that freedom.
The native theory must therefore explain the physical selection of the
constitutive structure and its normalization.  This selection requirement does not
exclude a future physical identity or a specifically GR-class
consistency theorem within independently justified premises.

Even a quadratic action/energy structure would close only part of the problem.  It could seed
linear dynamics and quadratic field stress, but $\beta=1$ depends on nonlinear self-source
information---cubic or higher-order structure in an action language, or its equivalent.  Thus
\begin{equation}
\boxed{
\text{linear native gravity completion}
\;\neq\;
\text{nonlinear native gravity completion}.}
\end{equation}

\subsection{Remaining native obligations}

At the scientific statuses used in this manuscript, the remaining obligations are:
\begin{enumerate}[leftmargin=*]
  \item whether the Lorentz-compatible vacuum is genuinely emergent/frameless rather than
        a fixed but locally undetectable background;
  \item the physical vacuum perturbation variables and native local response/evolution law;
  \item native selection of the long-range gravitational mode, including its gaplessness,
        rank/helicity, constraints or redundancies, and physical polarization count;
  \item the generic localized source functional and a consistent conserved source structure;
  \item universal gravitational coupling, its sign, source identity, and normalization as $G$;
  \item a universal finite-closure response theorem and any resulting native
        equivalence-principle statement;
  \item native interacting field energy/stress and the source/projection map to a
        gauge-consistent effective gravitational self-source, including a specific
        Weinberg--Witten applicability/evasion analysis and the nonlinear GR-class
        completion required for $\beta=1$;
  \item strong-field gravity, including horizon/interior structure, singularity behavior,
        mergers, ringdown, and compact-object response;
  \item a microscopic map, if one exists, between \cmotion{} budget variables and the
        effective metric geometry;
  \item the universal mass--radius bridge and accelerated-closure stability spectrum needed
        to determine any defensible acceleration coefficient.
\end{enumerate}

A complete microscopic action could in principle organize several of these structures at
once, but it is not logically required merely to formulate and test the first linear
vacuum-mode problem.  Conversely, the nonlinear self-source problem cannot be closed by an
arbitrary response operator alone; it needs a deeper conserved field-energy/stress structure.
These sharpened open gates are not publication blockers for the present weak-field companion.

None of these open problems is repaired by repeating standard GR mathematics in \cmotion{}
terminology.  Conversely, their existence does not erase the two useful results already
established within scope: the pure scalar/rate route fails, and an explicitly assumed massless
rank-2 completion is compatible with the required weak-field phenomenology.

\subsection{Conditional observational tests and later strong-field work}
\label{sec:conditional-tests}

The admitted pure massless helicity-2 completion has identifiable tests.
A reproducible additional radiative polarization, or a propagation law
incompatible with its assumed local-$c$ massless branch, would reject that
completion under its stated antecedents.  Locally measured propagation
must be distinguished from coordinate speeds and from source/emission
delays.  Such a failure would identify a failed branch or premise; it would
not by itself adjudicate every possible microscopic realization.

Possible nonlinear extensions include source saturation and regular
compact interiors.  These remain Conjectural mechanisms requiring a
specified nonlinear source/evolution law, conserved exchange, stable
solutions, and exterior matching.  Only after those steps could tidal
deformability, ringdown, or near-horizon lensing become quantitative
discriminators of a particular model.  Regularity alone supplies neither
a definite Love number nor a unique observational signature.  None of
these strong-field hypotheses is used to establish the present weak-field
results.

%======================================================================
\section{Conclusion}
\label{sec:conclusion}

Gravity is a nontrivial test of \cmotion{} precisely because Newtonian attraction and
clock-rate change are not enough.  The simplest candidate---one scalar gravitational load
that changes only local propagation/rate---captures part of the weak-field intuition but
fails the relativistic light-bending requirement.  In the tested class it gives
$\gamma=0$ and therefore half the observed leading deflection.  The missing piece is a
spatial or directional response.

A massless symmetric rank-2 field supplies that structure in a controlled effective
completion.  Within the selected vacuum route, with the stated source,
universal-coupling, common
probe-response, and linear-gauge antecedents, and no relevant extra scalar
admixture, it reproduces $\gamma=1$, full leading bending and Shapiro delay,
and the two luminal tensor gravity-wave polarizations.  Under the stated
localized slow-source and far-zone conditions, quadrupole radiation is
leading; with the additional specified GR-class nonlinear completion,
$\beta=1$ and the standard first post-Newtonian behavior follow.  These positive results show compatibility with the
observed weak-field gravity sector.  They do not amount to a first-principles derivation
of general relativity from \ElEmEr{}.

The candidate interpretation suggests a physical picture for what that
effective structure might mean.  A gravitational source may reorganize the local
field/stress state so that open propagation and recurrent return conditions vary with
position and source structure.  Free fall can then be read as the natural collective
response to that organization, support as resistance to it, clock shifts as changes of
physical recurrence, and tides as differential response across an extended body.  The
usual metric/geodesic description is the precise weak-field geometry of the resulting
motion.  This is a candidate \cmotion{} reading, not the missing microscopic law.

The closure-projection tests also preserve a useful boundary.  A scalar
steering model with an assumed universal half-response can produce a
pointwise geometric ray/matter coefficient ratio of two, but this alone
is not an integrated deflection prediction.  Raw isotropic averaging and the tested
strain, tidal, and retarded mechanisms do not derive that response.
The successful Conditional tensor comparison is therefore not a native
uniqueness proof.

The remaining question is sharper than at the start: what microscopic
\cmotion{} dynamics could realize the massless, universally coupled
tensor sector tested here, or another completion satisfying the same
observational requirements?  The open bridge is now localized: identify the
physical vacuum perturbation variables and their local response law; derive the generic
localized source and its universal coupling; establish how finite recurrent closures respond;
and obtain a self-consistent native energy/source law together with a
microscopic-to-effective gravitational stress/self-source projection so that
nonlinear completion can begin.  A full microscopic action is one possible route through several of
these stages, but the linear mode problem can be posed before that if a native constitutive
response law is found.  The separate acceleration-instability idea additionally requires a
solved accelerated closure and its stability spectrum.  Until those are known, the effective
geometric description remains the successful gravitational language while the field-first
interpretation remains an explicitly ongoing microscopic research problem.

%======================================================================
\newpage
\appendix
\section{Formal claim-status and evidential-role register}
\label{app:status}

The body of the paper is organized by the physical argument rather than by status
bookkeeping.  This appendix records the claim status of the
results used in this manuscript.  The scientific-status column uses the synchronized
C-Motion vocabulary.  Conditional means a conclusion follows from an explicit
premise or model class not itself derived; Structural means an exact or banked
organizational/accounting result without a supplied microscopic law;
Model-derived means rigorously derived within a specified effective/analogue
model, never native Derived.  Assumed denotes a provisional active premise;
Conjectural a motivated but unestablished mechanism or interpretation; Open an
unresolved question, with Open (sharp) marking an isolated missing bridge rather
than greater confidence.  Negative excludes only the stated identification,
mechanism or branch.  Publication deferral leaves the scientific status intact.

The role column is separate: N1 is an inherited native premise/constraint route;
N2 a C-Motion-constrained effective bridge; C a supplied standard/control
calculation; O an unresolved physical bridge.  The scalar/rate Negative retains
its N1 provenance flag for the framework-motivated restriction being tested;
this does not promote its supplied scalar ansatz or PPN comparison to native
dynamics.  Its mathematical comparison remains a scoped control.
R is a local rename-risk warning, not a status or an evidential upgrade:
rephrasing an effective geometry as a physical ``budget'' does not derive it.
It appears only on that interpretation when treated as evidence and is
unrelated to the radius $R$ in Section~\ref{sec:acceleration}.

{\footnotesize
\setlength{\tabcolsep}{3pt}
\begin{longtable}{@{}P{3.8cm}P{2.8cm}P{1.5cm}P{5.3cm}@{}}
\toprule
\textbf{Claim / result} & \textbf{Scientific status} & \textbf{Role} & \textbf{Exact scope / dependency} \\
\midrule
\endhead
Primitive local-$c$ restriction on an open gravity influence & \statusAssumed{} framework premise & \NOne{} & Openly propagated generated change respects local $c$; does not identify gravity source, rank, or coupling. \\
Complete-field source $\to$ local field/stress reorganization $\to$ response architecture & \statusConjectural{} & \NTwo{} & Common architecture only; Structural bookkeeping is carried in the scope, while the proposed sector realization is Conjectural. It does not derive the gravity source, mediator, $G$, or universal coupling. \\
Frameless Lorentz-invariant vacuum branch & \statusOpen{} & \OpenBridge{} & F selects the proposed native realization; effective fixed-background Fierz--Pauli mathematics neither proves nor independently requires it. \\
Scalar/rate-only loading gives $\gamma=0$ and half bending & \statusNegative{} & \NOne{} & Isotropic weak-field optical $v=cN$ (not primitive local $c$), $U=-\Phi_N/c^2$, no independent spatial response. \\
Viable gravity needs more than pure scalar rate loading & \statusStructural{} & \NTwo{} & Restricted rate/amount-only model fails; scalar derivative/projection alternatives remain open; rank 2 not uniquely selected. \\
Amount-only source-capacity restriction & \statusNegative{} & \NTwo{} & Rotationally scalar, linear, first-order local amount-only source with no independent angular-stress or reconstruction data cannot distinguish equal-$U_0$, zero-flux configurations with different $Q_{ij}$; not a general scalar no-go or rank-2 derivation. \\
Observed $G$-strength gravity is not the universal elementary closure binder & \statusNegative{} & \NTwo{} & Excludes ordinary gravity as universal elementary binder; actual binder remains open. \\
Massless symmetric rank-2/helicity-2 mediator M & \statusConditional{} & \OpenBridge{} / \Control{} & Not derived by local $c$, scalar failure, or common response architecture. \\
Native perturbation variables and local vacuum-response operator & \statusOpen{} & \OpenBridge{} & Local $c$ constrains characteristics only after a physical mode exists; component count, gap, rank, helicity and constraints remain undetermined natively. \\
Symmetric conserved total source S & \statusConditional{} & \OpenBridge{} / \Control{} & Effective rank-2 control requires a conserved symmetric source; a local native source functional/conservation law remains missing. An action is one route, not the only logically possible route. \\
Universal coupling and $G$ normalization U & \statusAssumed{} & \OpenBridge{} & Universality is used here as an Assumed antecedent; its native derivation, sign, source identity, and normalization remain Open. \\
Common effective probe-metric response P & \statusConditional{} & \Control{} & Standard comparison premise for light/clock/matter response; a native closure response theorem remains Open. \\
Universal finite-closure response & \statusOpen{} & \OpenBridge{} & Universal sourcing does not by itself prove composition-independent center response of differently structured recurrent closures. \\
Linear Fierz--Pauli class L & \statusConditional{} & \Control{} & Standard local two-derivative massless rank-2 gauge completion. \\
Closure-projector cycle-average and isotropic $2/3$ identity &
\statusStructural{} & \Control{} &
Exact within the assumed normalized cycle measure and projector;
not a native closure-response law. \\
Scalar steering plus universal half-response lock &
\statusConditional{} & \OpenBridge{} / \Control{} &
Pointwise geometric ray/matter ratio only; transverse ray, half-lock and
observable map supplied; no integrated bending, complete PPN, clock,
source or WEP derivation. \\
Tested strain/tidal/retarded half-lock rescue mechanisms &
\statusNegative{} & \Control{} &
Scoped failures derived in Appendix~\ref{app:projection-diagnostics};
not exclusions of every dynamical scalar or closure-memory model. \\
Native physical closure-centre projection and memory law &
\statusOpen{} & \OpenBridge{} &
No universal native averaging measure, non-raw projection, or closure-memory
law is supplied; recurrence does not establish a universal half-lock. \\
Rest-energy/inertial-mass bridge &
\statusConditional{} & \NTwo{} &
Complete isolated recurrent solution, conserved total energy--momentum,
rest frame and effective Lorentz dispersion; no gravitational equivalence
result. \\
Restricted extra scalar/vector source-class observations & \statusConditional{} & \Control{} & Not general no-go theorems and not part of the central reader-facing spine. \\
$\gamma=1$ & \statusConditional{} & \Control{} & Selected route F+M+S+U+P+L, with no relevant extra scalar; fixed-background effective control does not independently require F. \\
Full leading light bending and Shapiro delay & \statusConditional{} & \Control{} & Downstream of $\gamma=1$. \\
Luminal gravitational-wave propagation & \statusConditional{} & \NOne{} + \Control{} & Local $c$ constrains propagation; physical rank-2 wave additionally requires M/L. \\
Two TT/helicity-2 gravitational-wave polarizations & \statusConditional{} & \Control{} & M+L. \\
Quadrupole-leading radiation & \statusConditional{} & \Control{} & Isolated localized weak-field/slow-motion source, far-zone conditions, conserved source S, no extra radiative charges, and pure helicity-2 sector. \\
$\beta=1$ & \statusConditional{} & \Control{} & Requires specifically GR-class nonlinear N, correct constraints and same coupling, plus P; generic nonlinear self-interaction and quadratic seed are insufficient. \\
Free-transport budget/flux moment identities & \statusConditional{} & \NTwo{} & Exact within supplied source-free $\rho$ transport; not interacting field-energy conservation or gravitational source derivation. \\
Microscopic/effective gravitational stress bridge & \statusOpen{} & \OpenBridge{} & Native interacting energy/source, gauge-consistent effective self-source projection and Weinberg--Witten compatibility unproved; native local stress is not forbidden. \\
Native field-energy/stress/self-source structure & \statusOpen{} & \OpenBridge{} & Required for nonlinear gravitational self-coupling; an action/Hamiltonian/equivalent consistent route remains unclosed. \\
Standard 1PN perihelion consequence & \statusConditional{} & \Control{} & Follows after $\gamma=\beta=1$; GR/PPN-seeded numerics are controls only. \\
Native equivalence principle & \statusOpen{} & \OpenBridge{} & Requires universal coupling, source-to-response map, finite-closure response, and control of internal-structure corrections. \\
Budget geometry & \statusConjectural{} & \RenameRisk{} if treated as evidence & Interpretive provenance only; no native budget-to-metric derivation is claimed. \\
Open/massless-sector reading of the admitted gravity mode & \statusConjectural{} & \OpenBridge{} & Interpretive provenance only. The conditional massless mediator fits the open-propagation taxonomy; taxonomy does not derive M or establish a universal open$\leftrightarrow$massless equivalence. \\
Free-fall/support/tidal physical reading & \statusConjectural{} with native EP \statusOpen{} & \OpenBridge{} & Interpretive provenance only. Uniform-response and differential-response intuition does not supply a finite-closure averaging theorem or native WEP derivation; native EP remains Open. \\
Universal local-response / closure-cycle projection realization &
\statusConjectural{} &
\OpenBridge{} &
Candidate microscopic realization in which the same local gravitational response acts on
open and recurrent ElEmEr propagation, while recurrence changes only the net organized
response through cycle projection/cancellation.  No universal physical averaging law,
physical projection factor, acceleration dependence, or native equivalence-principle result is
derived. \\

Full nonlinear GR and strong-field completion & \statusOpen{} & \OpenBridge{} & Publication-deferred; scientifically Open. Requires N plus source/action/coupling completion. \\
Idealized two-state full-sweep branch mixing & \statusModelDerived{} & \Control{} & Supplied constant-gap Hamiltonian, phase rule and prescribed effective acceleration give $P_{\rm LZ}=e^{-\pi a_*/a_{\rm prop}}$; not physical closure destruction or a universal threshold. \\
Proper-acceleration closure breakdown $\sim C_{\rm dyn}c^2/R$ & \statusConjectural{} & \OpenBridge{} & Dimensional candidate scale only, not an instability theorem; requires an accelerated-closure stability calculation. \\
$a_c\sim mc^3/\hbar$ & \statusConditional{} & \OpenBridge{} & Conditional algebraic form of the Conjectural physical scale only: additionally requires $R=C_R\hbar/(mc)$; $C_a=C_{\rm dyn}/C_R$ remains unknown. \\
\bottomrule
\end{longtable}
}

\subsection*{Restricted source-class observations}
Two bounded observations concern simple additional long-range channels.  Given F, S, U, linear weak-field superposition, and no independent gravity
charge, a trace-coupled scalar is not sourced by traceless radiation, while a direct
$T^{00}$ scalar source selects a time direction and conflicts with F.  For a simple
vector source, $J^\mu=\mathcal T^{\mu\nu}n_\nu$ requires a preferred vector $n_\nu$,
whereas $J^\mu=\partial_\nu\mathcal T^{\mu\nu}$ vanishes by conservation.  These remain
conditional observations inside that restricted source-construction class.  They are not
general no-go theorems for scalar--tensor, vector--tensor, nonlocal, or additional-field
completions and are intentionally kept out of the main narrative.

\section{Definitions and derivations of the projection diagnostics}
\label{app:projection-diagnostics}

This appendix defines and derives the three summarized tests in
Section~\ref{sec:closure-projection}.  All conclusions are confined to the
specified projector, cycle measure and expansions.  They are analytical
diagnostics, not independent validation of physical closure response.
No native force law, elastic energy, closure stability theorem or universal
centre-response map is derived.

\subsection{Minimum differential steering strain}
\label{app:projection-strain}

Let $\mathbf g_0\ne0$ be a prescribed uniform
\emph{acceleration-valued diagnostic drive}.  The diagnostic forcing is
$\mathbf f(\mathbf n)=P(\mathbf n)\mathbf g_0$, with
$P(\mathbf n)=I-\mathbf n\mathbf n^{\mathsf T}$.
This definition supplies units consistently; it does not identify
$\mathbf f$ with a native force law or with $\mathrm d\mathbf n/\mathrm ds$.
For unit directions sampled uniformly in a plane with unit normal
$\mathbf N$, define
\begin{equation}
 C_N=\frac12(I-\mathbf N\mathbf N^{\mathsf T}),\qquad
 \mu=\mathbf N\cdot\widehat{\mathbf g}_0,\qquad
 \langle\mathbf f\rangle=(I-C_N)\mathbf g_0 .
 \label{eq:diag-plane-drive}
\end{equation}
Choose, solely as a tested orientation criterion, the residual variance
\begin{equation}
 \mathcal V_N
 :=\left\langle|\mathbf f-\langle\mathbf f\rangle|^2\right\rangle
 =\langle|\mathbf f|^2\rangle-|\langle\mathbf f\rangle|^2 .
 \label{eq:diag-variance-definition}
\end{equation}
Since $P^{\mathsf T}P=P$, direct angular averaging gives
\begin{align}
 \langle|\mathbf f|^2\rangle
 &=|\mathbf g_0|^2-\mathbf g_0^{\mathsf T}C_N\mathbf g_0
   =\frac{|\mathbf g_0|^2}{2}(1+\mu^2),\\
 |\langle\mathbf f\rangle|^2
 &=\frac{|\mathbf g_0|^2}{4}(1+3\mu^2).
\end{align}
Consequently
\begin{equation}
 \boxed{\mathcal V_N=\frac{|\mathbf g_0|^2}{4}(1-\mu^2).}
 \label{eq:diag-variance-result}
\end{equation}
Minimizing this particular variance favors $|\mu|=1$, where the
longitudinal response coefficient is one and $\mathcal V_N=0$.
The desired half-response orientation $\mu=0$ instead has variance
$|\mathbf g_0|^2/4$.  Thus the proposed minimum-variance criterion
does not select half-lock.  The criterion is not a derived physical
energy or stability functional, so this calculation excludes that
criterion as the claimed rescue, not every possible orientation law.



\subsection{Symmetric first-gradient tidal cancellation}
\label{app:projection-tidal}

Take orthonormal $\mathbf e_1,\mathbf e_2$, radius $R>0$, and
$\mathbf N=\mathbf e_1\times\mathbf e_2$.  Use the circular diagnostic
\begin{align}
 \mathbf y(\theta)&=R(\cos\theta\,\mathbf e_1+\sin\theta\,\mathbf e_2),\\
 \mathbf t(\theta)&=-\sin\theta\,\mathbf e_1+\cos\theta\,\mathbf e_2,
 \qquad \mathbf n_\pm(\theta)=\pm\mathbf t(\theta),\\
 P(\theta)&=I-\mathbf t(\theta)\mathbf t(\theta)^{\mathsf T}.
 \label{eq:diag-ring}
\end{align}
Here $\mathbf y$ is displacement from the closure centre; it is not the
dimensionless response $\mathbf r$ of
Eq.~\eqref{eq:projection-diagnostic}.  At each position the two directions
have equal counterflow weights.  The spatial cycle average also uses a
positive normalized weight $w$ with
$w(\theta+\pi)=w(\theta)$; opposite propagation directions alone would
not ensure this opposite-point symmetry.

At fixed evolution parameter, let a smooth prescribed external drive
vary on a scale $L\gg R$.  Retaining only first spatial order gives
\begin{equation}
 \mathbf g(\mathbf x_c+\mathbf y)
 =\mathbf g_0+\mathsf K\mathbf y+\cdots,\qquad
 K_{ij}=\partial_jg_i\big|_{\mathbf x_c}.
 \label{eq:diag-gradient-expansion}
\end{equation}
$\mathsf K$ is constant across this first-order calculation and has
units of acceleration per length; the omitted terms involve higher
spatial derivatives.  The identities
\begin{equation}
 \mathbf y(\theta+\pi)=-\mathbf y(\theta),\qquad
 P(\theta+\pi)=P(\theta)
\end{equation}
make $P(\theta)\mathsf K\mathbf y(\theta)$ odd under a half-turn.
Pairing opposite points in the weighted integral therefore gives
\begin{equation}
 \boxed{\left\langle P(\theta)\mathsf K\mathbf y(\theta)
 \right\rangle_w=0.}
 \label{eq:diag-gradient-zero}
\end{equation}
The first-gradient collective forcing cancels within this symmetric
diagnostic; the local differential forcing need not vanish and can
remain as internal tidal strain.  This calculation does not cancel
higher spatial orders or asymmetric cycle weights, nor does it prove
how an actual closure converts local forcing into centre motion.



\subsection{Retarded first-hit timing and odd memory}
\label{app:projection-retarded}

Use the same ring and opposite-point weighting.  Scale the compatible
evolution label $s$ so that $\cmc\,\Delta s$ is a propagation distance;
it is not a universal observer-clock reading.  Model an incident change
with approximately planar arrival surfaces across the ring, whose
source separation $L$ is much larger than $R$.  For a fixed unit
arrival axis $\hat{\mathbf a}$, the relative arrival model is
\begin{equation}
 s_{\rm ret}(\mathbf y)
 =s_0-\frac{\hat{\mathbf a}\cdot\mathbf y}{\cmc}
   +O\!\left(\frac{R^2}{L\cmc}\right).
 \label{eq:diag-arrival-time}
\end{equation}
The sign depends on the chosen direction of the arrival axis and does
not affect the symmetry cancellation.  Let $\mathbf g$ describe the
prescribed driving profile before this relative arrival shift, with
variation time $T$ satisfying $R/(\cmc T)\ll1$.  Taylor expansion to
first order in spatial displacement and arrival delay yields
\begin{equation}
 \mathbf g(\mathbf x_c+\mathbf y,s_{\rm ret})
 =\mathbf g_0+\mathsf K\mathbf y
  -\frac{\hat{\mathbf a}\cdot\mathbf y}{\cmc}\,
    \partial_s\mathbf g_0+\cdots .
 \label{eq:diag-arrival-expansion}
\end{equation}
All coefficients are evaluated at $(\mathbf x_c,s_0)$.
The ellipsis denotes omitted higher spatial/time Taylor terms, front-curvature
corrections and mixed higher orders.  The arrival shift must not be
applied again to driving data already evaluated with their full
retarded spatial dependence.

The changing-field diagnostic forcing is
\begin{equation}
 \mathbf f_{\rm ret}(\theta)
 =-\frac{\hat{\mathbf a}\cdot\mathbf y(\theta)}{\cmc}
    P(\theta)\partial_s\mathbf g_0,\qquad
 \boxed{\langle\mathbf f_{\rm ret}\rangle_w=0.}
 \label{eq:diag-retarded-zero}
\end{equation}
It vanishes pointwise for a settled static drive,
$\partial_s\mathbf g_0=0$.  For a changing drive its average still
vanishes because the scalar arrival delay is odd under a half-turn,
$P$ is even, and $w$ has the stated symmetry.  Internal differential
forcing can remain even when the collective average is zero.

The separate first-order memory test uses a dimensionless parameter
$\epsilon$ and an explicitly odd leading weight
$\chi(\theta+\pi)=-\chi(\theta)$:
\begin{equation}
 \mathcal R_{\rm mem}
 =\left\langle[1+\epsilon\chi(\theta)+O(\epsilon^2)]P(\theta)
   \right\rangle_w,\qquad
 \boxed{\langle\chi P\rangle_w=0.}
 \label{eq:diag-odd-memory-zero}
\end{equation}
Odd arrival weighting thus leaves the response tensor for the
\emph{uniform base drive} unchanged at first order.  Even memory,
broken cycle symmetry, nonlinear response and mixed memory--gradient
terms such as $\epsilon\langle\chi P\mathsf K\mathbf y\rangle_w$
are outside this cancellation and may survive.  None is fixed by the
timing model alone.  The result excludes the specified leading
first-hit/odd-memory rescue, not every retarded closure-response law.




\newpage
\printbibliography

\end{document}
